\icode{N25L-EMI-MU-L3-261004-002}
\nceeID[4]{2025山东}{11}
（多选）
如图所示，两条相同的半圆弧形光滑金属导轨固定在水平桌面上，其所在平面竖直且平行，导轨最高点到水平桌面的距离等于半径，最低点的连线\(OO^\prime\)与导轨所在竖直面垂直。空间充满竖直向下的匀强磁场（图中未画出），导轨左端由导线连接。现将具有一定质量和电阻的金属棒\(MN\)平行\(OO^\prime\)放置在导轨图示位置，由静止释放。\(MN\)运动过程中始终平行于\(OO^\prime\)且与两导轨接触良好，不考虑自感影响，下列说法正确的是\paren[]

\begin{tikzpicture}[font=\scriptsize]
	\draw[thick] (0,2) arc(-180:0:2) (0,2)--++(40:1)coordinate(a) arc(-180:0:2)coordinate(b);
	\draw[dashed] (0,0)--(4,0)--++(40:1)--++(-4,0)--cycle (a)--++(0,-2) (b)--++(0,-2) (0,0)--(0,2) (4,0)--++(0,2) (2,0)node[below left]{\(O\)}--++(40:1)node[above right]{\(O^\prime\)};
	\draw[very thick,red] ($(2,2)+(-150:2)+(-140:.1)$)node[below]{\(M\)}--++(40:1.2)node[above right]{\(N\)};
\end{tikzpicture}
\begin{choices}
	\item \(MN\)最终一定静止于\(OO^\prime\)位置
	\item \(MN\)运动过程中安培力始终做负功
	\item 从释放到第一次到达\(OO^\prime\)位置过程中，\(MN\)的速率一直在增大
	\item 从释放到第一次到达\(OO^\prime\)位置过程中，\(MN\)中电流方向由\(M\)到\(N\)
\end{choices}
\begin{answer}
	ABD\qquad 提示：
\end{answer}